\clearpage
\section{Technician: how the methods work}
\label{sec:methods}

\subsection{One idea behind the methods}

For a batch of readings $x_1,\ldots,x_n$, let $\bar{x}$ be the mean, $s$ the
standard deviation, $\tilde{x}$ the median, and
\[
\mathrm{MAD}=\operatorname{median}_i |x_i-\tilde{x}|.
\]
For a time series, $x_t$ denotes the value at time $t$.

Most methods in this report can be described by three parts:

\begin{enumerate}
  \item a \textbf{score} that says how unusual a reading is;
  \item a \textbf{reference} that describes normal behaviour;
  \item a \textbf{threshold} that converts the score into a flag.
\end{enumerate}

Keeping these parts separate is useful. A method can have a sensible score
but a poor reference, or a good reference but an unsuitable threshold.

\subsection{Classical and robust outlier scores}

Grubbs' statistic for the most extreme observation is
\begin{equation}
G=\frac{\max_i |x_i-\bar{x}|}{s}.
\label{eq:grubbs}
\end{equation}
The null hypothesis is that all observations come from one normal
population. A critical value is calculated from the sample size and chosen
significance level \citep{grubbs1969}. This is appropriate when the sample
matches the assumptions. A sensor stream with a daily cycle and strong
temporal correlation does not look like one independent normal batch.

The modified z-score replaces the fragile mean and standard deviation by the
median and MAD:
\begin{equation}
M_i=\frac{0.6745(x_i-\tilde{x})}{\mathrm{MAD}}.
\label{eq:modz}
\end{equation}
Iglewicz and Hoaglin recommend flagging observations with
$|M_i|>3.5$ as a practical screening rule \citep{iglewicz1993}. The score is
robust to a few extreme values, but the behaviour still depends on which
observations are placed in the window and on the threshold used.

\subsection{Operational checks}

Automatic environmental QC normally uses several simple checks rather than
one statistical test. Table~\ref{tab:checks} gives the small set used in this
report.

\begin{table}[h]
\centering
\small
\caption{The main checks, what they ask, and their main limitation.}
\label{tab:checks}
\begin{tabular}{@{}p{2.7cm}p{5.2cm}p{6.0cm}@{}}
\toprule
Check & Basic question & Main limitation \\
\midrule
Range & Is the value physically or instrumentally possible? & Cannot see plausible but wrong values. \\
Step & Did the value change more than a chosen amount between samples? & Misses slow drift and small level shifts. \\
Persistence & Has the sensor stopped changing? & Exact equality misses a stuck sensor with small last-digit jitter. \\
Internal consistency & Do related quantities at one station obey physical relationships? & Shows a contradiction but not necessarily which sensor is wrong. \\
Spatial consistency & Does the station agree with nearby stations after normal offsets are allowed for? & Can flag real local weather; depends on having suitable neighbours. \\
Outlier score & Is the reading unusual relative to a batch or moving window? & Depends strongly on the reference window and its assumptions. \\
Forecast & Is the reading far from what recent history predicted? & One-step forecasts often absorb slow drift and persistent faults. \\
\bottomrule
\end{tabular}
\end{table}

The checks should be read as different kinds of evidence. A range check uses
instrument knowledge. A step or persistence check uses the recent behaviour
of the same series. Internal consistency uses physics. Spatial consistency
uses redundancy across stations. Forecasting uses a fitted model of the
past. Their weaknesses are different, which is exactly why a battery can be
stronger than one check.

\subsection{Persistence as a useful example}

A stuck sensor can be tested in two ways. The simple implementation looks
for an exact run of repeated values:
\begin{equation}
\text{flag if }x_t\text{ belongs to a run of at least }n_{\mathrm{run}}
\text{ identical values.}
\end{equation}
A more robust version looks for a collapse in variability:
\begin{equation}
\text{flag window }W_t\iff \operatorname{sd}\{x_u:u\in W_t\}<s_{\min}.
\end{equation}
The second form is closer to the operational idea in \citet{shafer2000}.
It catches both a perfectly frozen value and a sensor that remains stuck but
jitters slightly in the last digit. This small example illustrates a wider
lesson: the exact definition of a check matters as much as its name.

\subsection{How a QC system should use the verdicts}

A useful QC pipeline should keep the original measurement and attach a
verdict or quality code rather than silently overwrite the data. A failed
range check may make a value unusable for later calculations, but the value
and the reason for rejecting it should remain auditable. Other checks usually
produce flags for later review because disagreement may represent either a
fault or a real event \citep{shafer2000,campbell2013}.

The order can also matter. Cheap checks can run before expensive spatial or
model-based checks, and clearly invalid values should not be used as evidence
for later checks. The important principle is transparency: the user should
be able to tell which check objected and why.


\subsection{The checks in plain language}

The table above is useful as a summary, but the checks are easier to explain
with one simple example each. I keep the descriptions here deliberately
practical, because these are the questions I would use when explaining the
system to an operator or to a reviewer.

\paragraph{Range check.}
Suppose an air-temperature sensor reports $-999$\,\textcelsius. The number
is not merely unusual; it is outside the operating range of the sensor and
outside any realistic environmental range. A range check compares the value
with a lower and upper limit and flags it immediately. This check is cheap
and easy to defend. Its weakness is that most real faults remain inside
plausible limits. A sensor that is wrong by 2\,K may still pass perfectly.

\paragraph{Step check.}
A step check asks whether two consecutive readings differ by more than a
chosen limit. For five-minute temperature data, a sudden jump of several
kelvin may be suspicious. This is good for spikes and abrupt failures. It is
not good for slow problems. A sensor that drifts by a tiny amount at every
sample can become badly wrong while never making one large step.

\paragraph{Persistence check.}
A persistence check looks for a sensor that has stopped responding. The
simplest version searches for exactly repeated values. That works if the
instrument sends the same number again and again. A more realistic faulty
sensor may still move in its last digit, so a low-variability version is
safer: instead of asking whether all values are identical, it asks whether
the variation inside a window has become abnormally small.

\paragraph{Internal-consistency check.}
Some quantities are related by physics. Dew-point temperature should not
normally exceed air temperature, and daily minimum temperature should not be
higher than daily maximum temperature. When two related quantities break a
physical relationship, at least one reading deserves attention. The limit
of the check is that it can show a contradiction without always identifying
which sensor caused it.

\paragraph{Spatial-consistency check.}
Nearby stations provide an external reference. If one station suddenly moves
away from several healthy neighbours while the neighbours continue to agree,
the station is suspicious. This is particularly useful for level shifts,
stuck sensors and drift. However, local weather can be real. A thunderstorm
may affect one station and not another, so disagreement is evidence for
review, not proof that a sensor is wrong.

\paragraph{Outlier-score check.}
An outlier score asks how far a reading lies from the centre of a reference
set. Grubbs uses the mean and standard deviation; the modified z-score uses
the median and MAD. These scores are useful only when the reference set makes
sense. A whole month of temperature contains a daily cycle, while a very
short window may be noisy. The main design question is therefore not only
``which score?'' but also ``which observations define normal?''

\paragraph{Forecast check.}
A forecast method predicts the next reading from recent history and flags a
large prediction error. A sudden spike is difficult to predict and therefore
easy to flag. Slow drift is different: if the model keeps updating from the
recent past, it can follow the drifting sensor and stop being surprised. The
method is therefore useful evidence for sudden changes, not a complete
solution to every fault.

These seven checks can be grouped into three simple sources of evidence.
Range and internal-consistency checks use physical knowledge. Step,
persistence, outlier and forecast checks use the behaviour of the same time
series. Spatial checking uses other stations. A strong QC system combines
these sources rather than asking one source to do everything.

\subsection{A simple running example}

Consider one station reporting air temperature every five minutes. At
10:00 it reports 12.1\,\textcelsius, at 10:05 it reports
12.2\,\textcelsius, and at 10:10 it reports 18.5\,\textcelsius. A range check
may accept all three values because 18.5\,\textcelsius is physically
possible. A step check may flag 10:10 because the jump is too large. A
rolling outlier score may also flag it because it is far from the recent
window. If neighbouring stations remain around 12\,\textcelsius, the spatial
check adds independent evidence that the value is suspicious.

Now consider a different failure. The sensor begins at 12.0\,\textcelsius
and then reports 12.01, 12.00, 12.02 and 12.01 for several hours while the
real temperature changes normally. The range check passes it and the step
check sees nothing unusual. Exact repetition may also miss it because the
values are not identical. A low-variability persistence check can detect the
collapse in movement, and the spatial check can detect the growing
disagreement with neighbouring stations.

These two examples explain why the pipeline needs several checks. The first
fault is sudden and visible in the local series. The second fault remains
plausible and requires a different definition of persistence or an external
reference. The checks are therefore complementary rather than competing.

% =====================================================================
% 4. Experimentalist
% =====================================================================
